\documentclass[11pt]{article}
\usepackage[letterpaper,margin=1in]{geometry}
\usepackage[T1]{fontenc}
\usepackage[utf8]{inputenc}
\usepackage{lmodern,microtype}
\usepackage{amsmath,amssymb,mathtools,bm,slashed,mathrsfs}
\usepackage{xcolor,enumitem,booktabs,array,needspace}
\usepackage[most]{tcolorbox}
\usepackage{fancyhdr,titlesec,parskip,hyperref}
\definecolor{qftblue}{HTML}{245A8D}
\definecolor{qftgreen}{HTML}{2F6B4F}
\definecolor{qftgray}{HTML}{F4F6F8}
\hypersetup{colorlinks=true,linkcolor=qftblue,urlcolor=qftblue,pdfauthor={Nima Lashkari}}
\pagestyle{fancy}\fancyhf{}
\lhead{PHYS 662: Quantum Field Theory I}\cfoot{\thepage}
\setlength{\headheight}{14pt}
\titleformat{\section}{\Large\bfseries\color{qftblue}}{\thesection.}{0.6em}{}
\titleformat{\subsection}{\large\bfseries\color{qftgreen}}{\thesubsection}{0.6em}{}
\setlist[enumerate]{leftmargin=2em,itemsep=0.25em,topsep=0.35em}
\setlist[itemize]{leftmargin=2em,itemsep=0.2em,topsep=0.35em}
\newcommand{\dd}{\mathrm d}
\newcommand{\ii}{\mathrm i}
\newcommand{\eps}{\varepsilon}
\newcommand{\tr}{\operatorname{tr}}
\newcommand{\Tr}{\operatorname{Tr}}
\newcommand{\diag}{\operatorname{diag}}
\newcommand{\Lag}{\mathcal L}
\newcommand{\D}{\mathcal D}
\newcommand{\vac}{\lvert0\rangle}
\newcommand{\vbra}{\langle0\rvert}
\newcommand{\Tord}{\mathrm T}
\newcommand{\ket}[1]{\lvert#1\rangle}
\newcommand{\bra}[1]{\langle#1\rvert}
\newtcolorbox{keyidea}[1][]{enhanced,breakable,colback=qftgray,colframe=qftblue,
 boxrule=0.7pt,arc=1.5mm,title={#1},fonttitle=\bfseries,coltitle=white}
\newtcolorbox{remarkbox}[1][]{enhanced,breakable,colback=white,colframe=qftgreen,
 boxrule=0.6pt,arc=1.5mm,title={#1},fonttitle=\bfseries,coltitle=white}
\newcommand{\intp}{\int\!\frac{\dd^4p}{(2\pi)^4}}
\newcommand{\MSbar}{\overline{\mathrm{MS}}}
\newcommand{\module}[1]{\clearpage\section{#1}}
\newcommand{\braket}[2]{\langle#1\mid#2\rangle}
\newcommand{\dPi}{\dd\Pi}
\newcommand{\cM}{\mathcal M}
\newcommand{\cH}{\mathcal H}
\newcommand{\cS}{\mathcal S}
\newcommand{\Zp}{Z_{\rm pole}}
\newcommand{\Kallen}{\mathcal K}
\newcommand{\lecturefront}[4]{%
\rhead{Lecture #1}\hypersetup{pdftitle={PHYS 662 Lecture #1: #2}}
\begin{document}\pagenumbering{gobble}
\begin{titlepage}\centering\vspace*{1.05in}
{\scshape\Large PHYS 662 --- Quantum Field Theory I\par}
\vspace{0.35in}{\Huge\bfseries\color{qftblue}Lecture #1\par}
\vspace{0.2in}{\LARGE\bfseries #2\par}
\vspace{0.2in}{\large\color{qftgreen}#3\par}
\vspace{0.6in}{\large Nima Lashkari\par}{\large Fall 2026\par}
\vfill\begin{minipage}{0.86\textwidth}\small #4\par\medskip
Course page: \url{https://nimalashkari.com/quantum-field-theory-2026-part-i/}
\end{minipage}\vspace{0.5in}\end{titlepage}
\pagenumbering{roman}\setcounter{tocdepth}{1}\tableofcontents
\vfill\begin{remarkbox}[Place in the course]
Lectures 21--22 developed interacting correlators and scalar loops.
Lecture 23 defines the scattering states and the S-matrix.
Lecture 24 converts amplitudes into rates and cross sections.
Lecture 25 derives the reduction formula connecting these subjects.
Homework 8 covers this three-lecture sequence.
\end{remarkbox}
\small\textbf{Reading and conventions.} These notes use the course slides and
D.~Harlow, \href{https://www.mit.edu/~harlow/HarlowQFT1.pdf}{\textit{QFT I}},
Sections 8--9, with our covariant state normalization and mostly-minus metric.
Useful online companions are D.~Tong,
\href{https://www.damtp.cam.ac.uk/user/tong/qft/qfthtml/S3.html}{\textit{QFT}, Sections 3.6--3.7},
and D.~Buchholz and W.~Dybalski,
\href{https://arxiv.org/abs/math-ph/0509047}{\textit{Scattering in relativistic quantum field theory}},
Sections 1--3, for the scope of the asymptotic construction.
For textbook study, see Weinberg, \textit{The Quantum Theory of Fields},
Vol.~I, Chapter 3, and Peskin--Schroeder, \textit{An Introduction to QFT},
Chapters 4 and 7. The accompanying editorial memo records the online
sources consulted and the corrections to the slides.
\normalsize\clearpage\pagenumbering{arabic}}
\lecturefront{24}{Decay Rates and Cross Sections}
{From normalized amplitudes to measurable event rates}
{We derive the phase-space and flux factors rather than memorizing them.
Two-body examples connect the invariant amplitude to lifetimes and
scattering cross sections. Resonances and the optical theorem then explain
how these observables fit the analytic and unitary structure of the theory.}

\section{What an experiment measures}
The S-matrix of Lecture 23 gives a transition amplitude between normalized
wave packets. A plane-wave matrix element is a distribution,
\begin{equation}
 \langle f|S-1|i\rangle_c
 =\ii(2\pi)^4\delta^{(4)}(P_f-P_i)\cM_{fi},
 \qquad
 \langle p|q\rangle=2E_p(2\pi)^3\delta^{(3)}(\bm p-\bm q).
 \label{eq:amp}
\end{equation}
Squaring just \(\cM\) is not yet a probability. We must account for state
normalization, the observation time, and the number of available final
states. For a collision we must also divide by the incoming flux.

For two dilute, incoherent beams with local laboratory number densities
\(n_A,n_B\), define the instantaneous luminosity by
\begin{equation}
 \mathscr L(t)=\int\dd^3x\,n_A(\bm x,t)n_B(\bm x,t)v_{\rm M},
 \qquad \frac{\dd N_X}{\dd t}=\mathscr L(t)\,\sigma_X .
 \label{eq:luminosity}
\end{equation}
Narrow momentum distributions are understood; otherwise average the
momentum-dependent product inside the integral. For collinear beams,
\(v_{\rm M}=|\bm v_A-\bm v_B|\); Section 5 gives its general definition.
\(\mathscr L\) has units of inverse area per time, and \(\sigma_X\) is an
area. A cross section is not itself a dimensionless probability.

An alternative packet picture labels collisions by the transverse impact
parameter \(\bm b\). In the narrow-momentum limit, with packets that scan
all transverse separations and the usual incident normalization,
\begin{equation}
 \sigma_X=\int\dd^2b\,P_X(\bm b).
\end{equation}
The probability changes when a beam is focused; the extracted cross section
characterizes the interaction at the chosen incoming momenta. This
picture also explains why two incoming plane waves require a flux factor.

For a narrow unstable excitation, the corresponding observable is the
probability to decay per unit proper time. If its rest-frame width is
\(\Gamma\), exponential decay over its regime of validity gives
\begin{equation}
 N(t)=N(0)e^{-\Gamma t},\qquad
 \tau=\Gamma^{-1},\qquad t_{1/2}=\frac{\log2}{\Gamma}.
\end{equation}
\begin{keyidea}[The normalization plan]
Use a finite box and a finite observation time to derive a rate. Then
remove the regulators after combining the squared amplitude with final
state counting. One incoming particle gives a decay formula; two incoming
particles give a collision rate that must be divided by a flux.
\end{keyidea}
An exactly stable asymptotic particle cannot decay. The decay formulas
below describe a weakly unstable prepared excitation to leading order,
or a narrow resonance whose mass and width are well separated scales.

\module{Why a squared delta function becomes time and volume}
Let \(\mathcal T\) be a long observation time, distinct from the transition
operator \(T\). Regulate energy conservation by
\begin{equation}
 F_{\mathcal T}(\omega)=\int_{-\mathcal T/2}^{\mathcal T/2}\dd t\,
 e^{\ii\omega t}=\frac{2\sin(\omega\mathcal T/2)}{\omega}.
\end{equation}
Its squared modulus has the distributional limit
\begin{equation}
 \frac{|F_{\mathcal T}(\omega)|^2}{\mathcal T}
 \longrightarrow 2\pi\delta(\omega).
 \label{eq:golden}
\end{equation}
Indeed, \(\int\dd\omega\,|F_{\mathcal T}(\omega)|^2/(2\pi)=\mathcal T\),
and the normalized peak narrows around zero. The limit is taken after
integration against a smooth final-state density. This is the origin of
Fermi's golden rule, not a prescription for multiplying arbitrary distributions.

In a large periodic box of volume \(V\), a momentum eigenstate has norm
\(2E_pV\). Its unit-normalized representative is \(\ket{p}/\sqrt{2E_pV}\),
and the number of final momentum modes in \(\dd^3p\) is
\(V\dd^3p/(2\pi)^3\). For the spacetime delta function these regulators give
\begin{equation}
 \big[(2\pi)^4\delta^{(4)}(P_f-P_i)\big]^2
 \rightsquigarrow V\mathcal T\,(2\pi)^4\delta^{(4)}(P_f-P_i).
\end{equation}
The arrow means precisely the regulated limiting operation above.

For \(r\) unit-normalized incoming modes, division by all external norms,
summation over final modes, and division by \(\mathcal T\) yield
\begin{equation}
 \frac{\dd P_{i\to f}}{\dd t}
 =\frac{V^{1-r}}{\prod_{a=1}^{r}2E_a}\,
 \frac{|\cM_{fi}|^2}{\cS_f}\,\dd\Phi_n(P_i),
 \label{eq:master}
\end{equation}
\begin{equation}
 \dd\Phi_n(P)=(2\pi)^4\delta^{(4)}\!\left(P-\sum_jp_j\right)
 \prod_j\frac{\dd^3p_j}{(2\pi)^3\,2E_j},\qquad
 \cS_f=\prod_a n_a!.
\end{equation}
There is no symmetry factor for the prepared incoming beam particles.
\(\cS_f\) removes repeated counting when we integrate over all labeled
momenta of identical final particles. Spin sums will be specified separately.

\begin{remarkbox}[A useful volume check]
For one incoming particle the volume cancels. For two, the rate is
proportional to \(1/V\): two particles in a larger box meet less often.
The incident flux is also proportional to \(1/V\), so their ratio is a
finite cross section. In four dimensions \([\dd\Phi_n]=({\rm mass})^{2n-4}\).
\end{remarkbox}

\module{Decay widths, branching fractions, and time dilation}
Set \(r=1\) in \eqref{eq:master}. In the rest frame of an unstable scalar
excitation \(A\) of mass \(M_A\), the partial width is
\begin{equation}
 \boxed{\quad \dd\Gamma_{A\to X}
 =\frac{1}{2M_A}\frac{1}{\cS_X}
 \overline{|\cM_{XA}|^2}\,\dd\Phi_X(p_A).\quad}
 \label{eq:width}
\end{equation}
The bar denotes a sum over unobserved final spins and, for an unpolarized
initial ensemble, an average over its initial spin states. There is no
initial spin average for a specified polarization. For scalar particles
there are no spin sums or averages.

The total width and branching fractions are
\begin{equation}
 \Gamma_{\rm tot}=\sum_X\Gamma_X,\qquad
 B_X=\frac{\Gamma_X}{\Gamma_{\rm tot}},\qquad \sum_XB_X=1.
\end{equation}
Channels in this sum must be distinct physical outcomes. Counting both an
inclusive channel and its exclusive subdivisions would count events twice.

If the parent has laboratory energy \(E_A\), the same invariant phase-space
integral is divided by \(2E_A\) instead of \(2M_A\). Consequently
\begin{equation}
 \Gamma_{\rm lab}=\frac{M_A}{E_A}\Gamma_{\rm rest},\qquad
 \tau_{\rm lab}=\frac{E_A}{M_A}\tau_{\rm rest}.
 \label{eq:dilation}
\end{equation}
The Lorentz-invariant content is decay per unit proper time. For a packet
with appreciable energy spread one averages its survival probability over
the momentum distribution; a single dilation factor then need not suffice.

Consider, as a running example,
\begin{equation}
 \Lag_{\rm int}=-\frac{g}{2}A\chi^2,\qquad M_A>2m_\chi,
 \qquad \ii\cM(A\to\chi\chi)=-\ii g.
 \label{eq:decaymodel}
\end{equation}
Here \([g]={\rm mass}\). Assume additional interactions stabilize the
potential, the vacuum is chosen consistently, and \(g\) is small enough
that \(\Gamma_A\ll M_A\). The identical \(\chi\) particles require
\(\cS_f=2!\), even though the vertex rule already included the \(2!\)
Wick contractions. The two factorials have different origins.

\begin{keyidea}[Why this does not contradict Lecture 23]
In a leading decay calculation, \(A\) is treated as a stable external mode
of the zeroth-order theory to compute its small decay probability. At
nonzero coupling it is not an exact asymptotic state. Resumming its
self-energy produces the resonance description in Section 7. The stable
\(\chi\) particles remain legitimate asymptotic external states.
\end{keyidea}

\module{Two-body phase space from the delta functions}
Let \(P^2=s=W^2\), \(P^0>0\), and work in the rest frame \(P=(W,\bm0)\).
For final masses \(m_1,m_2\), the spatial delta function sets
\(\bm p_2=-\bm p_1\). Writing \(p=|\bm p_1|\) gives
\begin{equation}
 \dd\Phi_2=\frac{p^2\dd p\,\dd\Omega}{16\pi^2 E_1E_2}
 \delta(W-E_1-E_2).
\end{equation}
At its physical root \(p=p_*\),
\begin{equation}
 \frac{\dd(E_1+E_2)}{\dd p}=\frac{pW}{E_1E_2},\qquad
 p_* =\frac{\sqrt{\Kallen(s,m_1^2,m_2^2)}}{2\sqrt{s}},
\end{equation}
where the K\"all\'en kinematic polynomial is
\begin{equation}
 \Kallen(x,y,z)=x^2+y^2+z^2-2xy-2xz-2yz.
\end{equation}
It follows that
\begin{equation}
 \boxed{\ \dd\Phi_2=\frac{p_*}{16\pi^2 W}\dd\Omega,\qquad
 \int\dd\Phi_2=\frac{p_*}{4\pi W}
 =\frac{\sqrt{\Kallen(s,m_1^2,m_2^2)}}{8\pi s}.\ }
 \label{eq:two}
\end{equation}
These formulas apply for \(W\geq m_1+m_2\); below threshold the physical
phase space vanishes. A positive algebraic value of \(\Kallen\) on another
branch does not by itself define a physical final state.

For equal masses, \(p_*=W\beta/2\) with
\(\beta=\sqrt{1-4m^2/s}\), so \(\int\dd\Phi_2=\beta/(8\pi)\).
For an angle-independent decay amplitude,
\begin{equation}
 \Gamma_{A\to12}=\frac{p_*}{8\pi M_A^2\cS_f}|\cM|^2.
\end{equation}
The example \eqref{eq:decaymodel} therefore gives
\begin{equation}
 \boxed{\quad \Gamma_{A\to\chi\chi}
 =\frac{g^2}{32\pi M_A}\sqrt{1-\frac{4m_\chi^2}{M_A^2}}.\quad}
 \label{eq:examplewidth}
\end{equation}
The width has mass dimension one and vanishes at threshold. The square-root
suppression here comes from phase space; amplitudes with angular-momentum
barriers can vanish with additional powers of \(p_*\).

\begin{remarkbox}[A route to more final particles]
Phase space factorizes by introducing the invariant mass of a subgroup:
\[
 \dd\Phi_3(P;p_1,p_2,p_3)=\frac{\dd s_{12}}{2\pi}
 \dd\Phi_2(P;q,p_3)\,\dd\Phi_2(q;p_1,p_2),\quad q^2=s_{12}.
\]
The limits are \((m_1+m_2)^2\leq s_{12}\leq(W-m_3)^2\).
The final-state counting factor still belongs to the full physical channel.
\end{remarkbox}

\module{The invariant flux and the differential cross section}
For one incoming particle of each beam in a box, \eqref{eq:master} gives
\(\dd P/\dd t=|\cM|^2\dd\Phi_n/(4E_AE_BV\cS_f)\).
Dividing by the incident flux per target, \(v_{\rm M}/V\), yields
\begin{equation}
 \boxed{\quad \dd\sigma_{AB\to X}=
 \frac{1}{F}\frac{\overline{|\cM_{XA B}|^2}}{\cS_X}\dd\Phi_X,
 \qquad F=4\sqrt{(p_A\cdot p_B)^2-m_A^2m_B^2}.\quad}
 \label{eq:sigma}
\end{equation}
The square root and the phase-space measure are Lorentz invariant.
The quantity used in laboratory luminosities is
\begin{align}
 v_{\rm M}&=\frac{\sqrt{(p_A\cdot p_B)^2-m_A^2m_B^2}}{E_AE_B}\notag\\
 &=\sqrt{|\bm v_A-\bm v_B|^2-|\bm v_A\times\bm v_B|^2}.
 \label{eq:moller}
\end{align}
To verify the identity, write
\(p_A\cdot p_B=E_AE_B(1-\bm v_A\cdot\bm v_B)\) and
\(m_A^2m_B^2=E_A^2E_B^2(1-v_A^2)(1-v_B^2)\).
For collinear beams the cross product vanishes, recovering the familiar
relative speed in that frame.

The M\o ller factor \(v_{\rm M}\) is a flux factor, not the speed of one
particle measured in the rest frame of the other. It can exceed one for
oppositely directed relativistic beams; this does not violate causality.
For two massive particles the latter relative speed is instead
\(\sqrt{(p_A\cdot p_B)^2-m_A^2m_B^2}/(p_A\cdot p_B)\), which is below one.

Two useful special frames are
\begin{equation}
 F_{\rm CM}=4p_i\sqrt{s}=2\sqrt{\Kallen(s,m_A^2,m_B^2)},\qquad
 F_{B\,\mathrm{at\ rest}}=4m_B|\bm p_A|.
\end{equation}
In the center-of-mass frame \(p_i\) is either incoming momentum magnitude.
The cross section is invariant when expressed through invariant kinematics;
the separate laboratory density and velocity factors need not be.

\begin{keyidea}[Spin and counting conventions]
For an unpolarized two-particle beam, average over its incoming spin states
and sum over unobserved outgoing spins. Do not average over final spins.
Identical incoming particles in two prepared beams introduce no extra
\(1/2!\) into \eqref{eq:sigma}. Counting unordered pairs in a gas, such as
an event density \(n^2\langle\sigma v_{\rm M}\rangle/2\), is a separate step.
\end{keyidea}

\module{Two-to-two scattering and a complete scalar example}
Combine \eqref{eq:two} and \eqref{eq:sigma}. In the center-of-mass frame,
with initial and final momentum magnitudes \(p_i,p_f\),
\begin{equation}
 \boxed{\quad\frac{\dd\sigma}{\dd\Omega}
 =\frac{1}{64\pi^2s}\frac{p_f}{p_i}
 \frac{\overline{|\cM(s,t)|^2}}{\cS_f}.\quad}
 \label{eq:dsdo}
\end{equation}
For identical final particles this formula uses a full labeled angular
sphere with the factor \(1/\cS_f\); it is an event-counting convention.
One can instead integrate a nonredundant angular region without that
factor. A single-particle inclusive angular distribution is a different
observable and counts the detected particles, not just the events.

For \(A+B\to1+2\), set the incoming \(A\) direction as the polar axis.
The momentum-transfer invariant is
\begin{equation}
 t=m_A^2+m_1^2-2E_AE_1+2p_ip_f\cos\theta,\qquad
 \frac{\dd t}{\dd\cos\theta}=2p_ip_f.
\end{equation}
For an azimuthally symmetric amplitude, integration over the azimuth gives
\begin{equation}
 \frac{\dd\sigma}{\dd t}
 =\frac{\overline{|\cM(s,t)|^2}}
 {16\pi\Kallen(s,m_A^2,m_B^2)\cS_f},
 \quad
 t_\pm=m_A^2+m_1^2-2E_AE_1\pm2p_ip_f.
 \label{eq:dsdt}
\end{equation}
Use the same counting convention for the \(t\) range as for the sphere.

In the theory \(\Lag_{\rm int}=-\lambda\phi^4/4!\), all four particles
have mass \(M\), \(p_i=p_f\), and at tree level \(\cM=-\lambda\).
There are two identical final particles. Thus
\begin{equation}
 \frac{\dd\sigma}{\dd\Omega}=\frac{\lambda^2}{128\pi^2s},\qquad
 \boxed{\ \sigma_{\phi\phi\to\phi\phi}
 =\frac{\lambda^2}{32\pi s}+O(\lambda^3).\ }
 \label{eq:quarticsigma}
\end{equation}
The amplitude already includes exchange symmetry through its Feynman
rules; there is only one quartic contact vertex to draw. Adding another
``exchange contact diagram'' would double the amplitude incorrectly.

\begin{remarkbox}[Checks at threshold and at high energy]
For equal-mass elastic scattering, the vanishing final phase space and
incoming flux cancel, leaving a finite threshold limit of \(\sigma\).
The rate coefficient \(\sigma v_{\rm M}\) still tends to zero. At fixed
weak coupling and high energy the result scales as \(s^{-1}\), consistent
with the dimension of an area. When large logarithms matter, the loop
corrections and running discussed in Lecture 22 must also be included.
\end{remarkbox}

\module{Resonances: lifetime, propagator, and line shape}
An unstable excitation appears as a resonance in scattering of stable
particles. Continue its two-point function through the physical cut.
For a narrow, isolated resonance the nearby pole lies on the appropriate
unphysical sheet, approximately at
\begin{equation}
 s_{\rm pole}\simeq(M_R-\ii\Gamma_R/2)^2
 =M_R^2-\ii M_R\Gamma_R+O(\Gamma_R^2).
\end{equation}
Near the resonance and away from rapidly varying thresholds, its
propagator on the physical real-energy boundary has the form
\begin{equation}
 \widetilde G_R(s)\simeq
 \frac{\ii\mathcal Z_R}{s-M_R^2+\ii M_R\Gamma_R}.
 \label{eq:resprop}
\end{equation}
This is a local approximation to a function with a continuum spectral
measure, not an extra stable delta function in that measure.

The pole near positive rest energy is \(E\simeq M_R-\ii\Gamma_R/2\).
Fourier transformation for positive time gives the survival amplitude
and probability
\begin{equation}
 A(t)\simeq e^{-\ii M_Rt}e^{-\Gamma_Rt/2},\qquad
 |A(t)|^2\simeq e^{-\Gamma_Rt}.
 \label{eq:survival}
\end{equation}
This refers to a prepared rest excitation or a fixed-momentum pole
contribution, not a generic position-space propagator with packet spreading.
Exponential decay is an intermediate-time approximation; an exact
bounded-below spectrum also produces short-time and late-time deviations.

To connect with Lecture 22, write the denominator as
\(s-m^2-\Sigma_R(s)\), and choose \(M_R^2-m^2-\operatorname{Re}\Sigma_R(M_R^2)=0\).
Expanding its real part gives, to narrow-width accuracy,
\begin{equation}
 \mathcal Z_R=[1-\operatorname{Re}\Sigma_R'(M_R^2)]^{-1},\qquad
 M_R\Gamma_R=-\mathcal Z_R\operatorname{Im}\Sigma_R(M_R^2).
\end{equation}
This sign follows from our insertion convention \(-\ii\Sigma_R\).
\(\mathcal Z_R\) here is a resonance approximation, not the positive stable
pole residue used in LSZ.

For slowly varying production and decay numerators,
\begin{equation}
 |\cM(s)|^2\ \propto\ \frac{1}{(s-M_R^2)^2+M_R^2\Gamma_R^2},\qquad
 \frac{1}{(s-M_R^2)^2+M_R^2\Gamma_R^2}
 \longrightarrow\frac{\pi}{M_R\Gamma_R}\delta(s-M_R^2).
\end{equation}
The last relation holds under an integral when the narrow peak lies well
inside its range. In the energy variable \(\sqrt{s}\), the full width at
half maximum is approximately \(\Gamma_R\). Background interference and
nearby thresholds can spoil this simple line shape and its factorization.

\module{The optical theorem as a measurable normalization check}
For fixed incoming quantum numbers, Lecture 23 gave
\begin{equation}
 2\operatorname{Im}\cM_{ii}
 =\sum_X\frac1{\cS_X}\sum_{\rm final\ spins}
 \int\dd\Phi_X\,|\cM_{Xi}|^2.
\end{equation}
Comparing with \eqref{eq:sigma}, the right side is \(F\sigma_{\rm tot}\).
Therefore, for two incoming particles,
\begin{equation}
 \boxed{\quad\operatorname{Im}\cM_{ii}(s,t=0)
 =\frac{F}{2}\sigma_{\rm tot}(s)
 =2p_i\sqrt{s}\,\sigma_{\rm tot}(s).\quad}
 \label{eq:opticalrate}
\end{equation}
The incoming and outgoing states on the left are the same: this is the
\emph{forward elastic} amplitude. The right side sums all allowed final
channels, including inelastic ones. If initial spins are averaged, average
the diagonal forward amplitude over the same ensemble. This is not an
identity between a generic-angle amplitude and a total cross section.

Check the factors with \eqref{eq:quarticsigma}. For equal incoming masses,
\(F=2s\beta\), where \(\beta=\sqrt{1-4M^2/s}\). The optical theorem predicts
\begin{equation}
 \operatorname{Im}\cM_{\rm forward}^{\rm one\ loop}
 =\frac{\lambda^2\beta}{32\pi}.
\end{equation}
This is precisely the imaginary part of the physical \(s\)-channel bubble
found in Lecture 22. Equivalently,
\(2\operatorname{Im}\cM=\lambda^2\int\dd\Phi_2/2!\).
Dropping the identical-particle factor would fail this check by two.

\begin{keyidea}[A practical calculation in five steps]
Specify the incoming ensemble and the exclusive final channel.
Compute its on-shell \(\cM\) with the conventions of Lecture 21.
Sum unobserved final spins and average only the specified initial ensemble.
Integrate the allowed phase space, counting identical final states once.
Finally divide by \(2M_A\) for a rest-frame decay or by \(F\) for a collision.
\end{keyidea}
These rules separate dynamics, which is in \(\cM\), from kinematics and
normalization. They apply to scalar and fermion processes alike, provided
the amplitude and the asymptotic-state assumptions are appropriate.
The remaining theoretical step is to justify how the correlation
functions calculated by path integrals produce this physical amplitude.
That is the purpose of LSZ reduction in Lecture 25.
\end{document}
